% Layout of a research paper in a Nature Portfolio journal, rebuilt from the author guidance in
% writeups/templates/guidance/ and from the page layout of the stored exemplar papers in
% writeups/templates/papers/.
%
% This file reproduces a layout convention and it reproduces no publisher identity. It carries no
% journal name, no masthead, no logo and no volume line, because this document is not a paper of
% any journal. The running head names this repository.

% ---------------------------------------------------------------- page
\usepackage[a4paper,top=21mm,bottom=20mm,left=16mm,right=16mm,columnsep=6.5mm,
            headsep=5mm,footskip=11mm]{geometry}
\setlength{\parindent}{0pt}
\setlength{\parskip}{4pt plus 1pt minus 0.5pt}
\setlength{\emergencystretch}{2em}
\raggedbottom

% ---------------------------------------------------------------- fonts
\usepackage[T1]{fontenc}
\usepackage{newtxtext}
\usepackage{newtxmath}
\renewcommand{\sfdefault}{phv}

% ---------------------------------------------------------------- colour
\usepackage{xcolor}
% The accent is the teal of the exemplar page furniture. It marks a heading rule, a link and the
% running head, and it carries no meaning of its own.
\definecolor{paperaccent}{RGB}{0,104,118}
\definecolor{paperrule}{gray}{0.62}

% ---------------------------------------------------------------- headings
\usepackage{titlesec}
\titleformat{\section}
  {\sffamily\bfseries\fontsize{10.5}{12.5}\selectfont\color{black}}{}{0pt}{}
\titlespacing*{\section}{0pt}{12pt plus 3pt}{3pt}
\titleformat{\subsection}
  {\sffamily\bfseries\fontsize{9.2}{11}\selectfont\color{black}}{}{0pt}{}
\titlespacing*{\subsection}{0pt}{10pt plus 2pt}{2pt}
\titleformat{\subsubsection}
  {\sffamily\itshape\fontsize{9}{11}\selectfont\color{black}}{}{0pt}{}
\titlespacing*{\subsubsection}{0pt}{8pt plus 2pt}{2pt}

% ---------------------------------------------------------------- running head and foot
\usepackage{fancyhdr}
\usepackage{lastpage}
\providecommand{\paperrunninghead}{Article}
\providecommand{\paperrunningsource}{}
\providecommand{\paperrunningfoot}{}
\pagestyle{fancy}
\fancyhf{}
\renewcommand{\headrulewidth}{0.5pt}
\renewcommand{\footrulewidth}{0.5pt}
\renewcommand{\headrule}{\color{paperrule}\hrule height \headrulewidth}
\renewcommand{\footrule}{\color{paperrule}\hrule height \footrulewidth}
\fancyhead[L]{\sffamily\bfseries\fontsize{8}{10}\selectfont\paperrunninghead}
\fancyhead[R]{\sffamily\fontsize{8}{10}\selectfont\color{paperaccent}\paperrunningsource}
\fancyfoot[L]{\sffamily\fontsize{8}{10}\selectfont\color{paperaccent}\paperrunningfoot}
\fancyfoot[R]{\sffamily\bfseries\fontsize{9.5}{11}\selectfont\thepage}
\fancypagestyle{plain}{%
  \fancyhf{}%
  \renewcommand{\headrulewidth}{0pt}%
  \renewcommand{\footrulewidth}{0.5pt}%
  \fancyfoot[L]{\sffamily\fontsize{8}{10}\selectfont\color{paperaccent}\paperrunningfoot}%
  \fancyfoot[R]{\sffamily\bfseries\fontsize{9.5}{11}\selectfont\thepage}%
}

% ---------------------------------------------------------------- captions
% The exemplar caption opens with the item label, a vertical bar, a bold title and then the
% description. writeups/templates/figure_design.md rule 4 gives the form and quotes the model.
\usepackage{caption}
\DeclareCaptionLabelSeparator{bar}{ \textbar\ }
\captionsetup{
  font={small},
  labelfont={sf,bf},
  labelsep=bar,
  justification=justified,
  singlelinecheck=false,
  skip=4pt,
}
\captionsetup[figure]{name=Fig.}
\captionsetup[table]{name=Table,position=top}

% ---------------------------------------------------------------- tables
\usepackage{array}
\usepackage{booktabs}
\usepackage{tabularx}
\newcolumntype{L}{>{\raggedright\arraybackslash}X}
\renewcommand{\arraystretch}{1.15}

% ---------------------------------------------------------------- floats
\usepackage{float}
\setlength{\textfloatsep}{12pt plus 2pt minus 2pt}
\setlength{\floatsep}{12pt plus 2pt minus 2pt}
\setlength{\intextsep}{12pt plus 2pt minus 2pt}
\setlength{\dblfloatsep}{12pt plus 2pt minus 2pt}
\setlength{\dbltextfloatsep}{14pt plus 2pt minus 2pt}
\renewcommand{\topfraction}{0.9}
\renewcommand{\dbltopfraction}{0.9}
\renewcommand{\bottomfraction}{0.7}
\renewcommand{\textfraction}{0.08}
\renewcommand{\floatpagefraction}{0.75}
\renewcommand{\dblfloatpagefraction}{0.75}

% ---------------------------------------------------------------- references
% biblatex writes the separator of a page range and of a compressed citation range as an en dash.
% CLAUDE.md and docs/report_style.md forbid an en dash and an em dash in any text of this
% repository, so both separators become a single hyphen.
\AtBeginDocument{%
  \providecommand{\bibrangedash}{-}%
  \renewcommand{\bibrangedash}{-}%
  \providecommand{\citerangedash}{-}%
  \renewcommand{\citerangedash}{-}%
}
% The reference list is set small and tight, as in the exemplar papers.
\AtBeginDocument{%
  \providecommand*{\bibfont}{}%
  \renewcommand*{\bibfont}{\fontsize{7.6}{9}\selectfont}%
}

% ---------------------------------------------------------------- inline code
\let\mvatexttt\texttt
\renewcommand{\texttt}[1]{{\def\_{\textunderscore\discretionary{}{}{}}\small\mvatexttt{#1}}}
